\section{Introduction}
\label{sec:intro}

High-resolution aerial and maritime imagery is the regime where small-object
detection is hardest: targets occupy a handful of pixels, background dominates
the frame, and running a dense detector at full resolution wastes most of its
computation on regions with no object of interest \cite{ESOD}. A common
remedy is \emph{selective computation}: a lightweight early-stage network
scores which spatial regions merit full-resolution processing, and only those
regions proceed through the expensive backbone/neck/head
\cite{QueryDet,ESOD,CEASC}.

Existing selectors share an implicit assumption: a region containing a real
object will produce a sufficiently high foreground/objectness/query response
in the lightweight branch. This holds for medium and large objects but
becomes unreliable exactly where selective computation matters most --
genuinely tiny targets, whose downsampled feature footprint can collapse to
one or two activations indistinguishable from background clutter
\cite{SET}. The two possible selector errors are asymmetric: retaining an
extra background region only wastes compute, while dropping a true-object
region removes that object from every later stage. An objectness-only
selector's failure mode on tiny objects is therefore a silent, unrecoverable
false negative, not simply a lower AP.

We ask whether a second, independent evidence signal -- local
spectral/saliency response, motivated by the observation that tiny targets
retain residual high-frequency structure even when their semantic response is
weak \cite{SET} -- can recover objects an objectness-only selector would
drop, and whether directly supervising \emph{object-level coverage} (does at
least one candidate cell overlap each ground-truth box?) in addition to
per-cell classification yields a higher-recall selector under the same
routing rule.

We instantiate this idea as DES-ESOD, on ESOD's patch-routing skeleton
\cite{ESOD}: a dual-evidence priority head replaces ESOD's single
objectness head, trained with an explicit soft object-coverage loss. The
results in this paper use ESOD's own native fixed-threshold routing,
unchanged -- not a fixed top-$K$ cap. This holds the routing rule and threshold
fixed, but does not hold realized computation fixed because different score
distributions can retain different areas; a fixed-cardinality top-$K$ mode is
part of the wider design but is evaluated separately, outside this paper.
Because a higher-recall selector naturally
retains more image area, we pair it with an efficient decoupled head and
report retained area (Occupy) alongside accuracy so this
trade-off is never hidden. We define evaluation on two datasets outside ESOD's own
reported protocol -- UAVDT (vehicles) and TinyPersonV2/SeaPerson (extremely
small pedestrians) -- to test whether a dual-evidence, coverage-supervised
selector transfers across domains it was not tuned on.

\textbf{Contributions:} (1) an object-level soft-coverage loss that directly
supervises the event recall actually depends on, instead of a per-cell
classification proxy; (2) a lightweight semantic-spectral evidence fusion
that repurposes spectral saliency as a routing signal rather than a feature
enhancer; (3) an ISPP-based decoupled head that reduces the dense detection
head's parameter and arithmetic cost while preserving ESOD's prediction
contract; and (4) a cross-domain evaluation protocol on UAVDT and
TinyPersonV2, with final results retained as placeholders while experiments
are in progress.
